\documentclass[11pt,a4paper]{article}
\usepackage[utf8]{inputenc}
\usepackage[margin=2.2cm]{geometry}
\usepackage{amsmath,amssymb}
\usepackage{hyperref}
\usepackage{xcolor}
\usepackage{booktabs}
\usepackage{listings}
\usepackage{graphicx}
\usepackage{enumitem}
\usepackage{caption}
\usepackage{subcaption}
\usepackage{microtype}
\usepackage{parskip}
\usepackage{textcomp}

\hypersetup{colorlinks=true,linkcolor=blue!50!black,urlcolor=blue!60!black}

\lstset{basicstyle=\ttfamily\footnotesize,breaklines=true,
        frame=single,framerule=0pt,backgroundcolor=\color{black!4},
        columns=fullflexible,keepspaces=true,
        showstringspaces=false,xleftmargin=0pt}

\title{Julia v1.5 SMC$^2$-MPC \\
       6-horizon sweep with FILTER HMC OFF, all rejuvenation tricks ON \\
       \large $T \in \{14, 28, 42, 56, 70, 84\}$ days, RTX 5090, seed 42}
\author{}
\date{2026-05-09}

\begin{document}
\maketitle

\begin{abstract}
This report records a 6-horizon sweep of the FSA v1.5 closed-loop
SMC$^2$-MPC bench with the \emph{filter}-side HMC kernel disabled
(\texttt{--num-mcmc 0}) and all available rejuvenation tricks turned
on (Liu--West $\theta$-cloud shrinkage at $a=0.97$;
Silverman-bandwidth KDE resample with coupled Liu--West correction
at \texttt{--smooth-resample-bw 1.0}; multivariate Gaussian bridge
between rolling windows). The controller-side HMC (ChEES, in
\texttt{Control/GPUControlSMC.jl}) is unchanged but is run at a
\textbf{reduced ``fast-controller'' particle / MCMC configuration}:
\texttt{--ctrl-n-smc} dropped $2048 \to 1024$ and
\texttt{--ctrl-num-mcmc} dropped $16 \to 8$ relative to the
saturated default (per the technical guide's
``shorter-wall-time'' recipe). Per-tempering-level controller
diagnostics are recorded for the first time. All numerical entries
in this report are pulled directly from the run artefacts on disk;
the parameter-recovery claim in \S\ref{sec:conclusion} is the senior
engineer's reading of the posterior-trace figures.
\end{abstract}

\tableofcontents
\clearpage


% =====================================================================
\section{Purpose}
\label{sec:purpose}

This study is a \textbf{baseline data-gathering run}: how does the
v1.5 filter perform on parameter recovery when the filter HMC is
removed entirely and the cloud is kept alive solely by the
rejuvenation tricks (Liu--West, Silverman-KDE, Gaussian bridge)?

The motivation is the prior finding (Tier~0,
\texttt{TIER\_0\_FINDINGS.md} in the sibling sweep folder
\texttt{julia\_horizon\_sweep\_2026-05-09/}) that the filter HMC
accepts 0\,\% of proposals across 2880 logged tempering levels at the
saturated configuration. If the HMC is doing nothing, removing it
should have no effect on filter posterior quality, and the runtime
should drop. This sweep measures that.

The output is a \textbf{baseline} that future studies (with HMC
re-enabled and tuned, or with different rejuvenation knobs) can
compare against.


% =====================================================================
\section{What was run}
\label{sec:cli}

Per horizon, identical except \texttt{--T-days}:

\begin{lstlisting}[language=bash]
julia --project=. tools/bench_smc_full_mpc_fsa_gpu.jl \
    --T-days <T> --seed 42 \
    --num-mcmc 0 --hmc-step-size 0 --hmc-leapfrog 0 \
    --liu-west-a 0.97 \
    --smooth-resample-bw 1.0 \
    --gaussian-bridge true \
    --N-smc 512 --K-per-chain 1000 \
    --ctrl-n-smc 1024 --ctrl-num-mcmc 8 \
    --collect-ctrl-diagnostics true \
    --output-dir <SWEEP_ROOT>/T<T>d_seed42
\end{lstlisting}

Driven by the sibling launcher
\texttt{version\_1\_5\_Julia/tools/launchers/run\_julia\_horizon\_sweep\_no\_filter\_hmc.sh}.

\subsection{What each filter-side flag does}

\begin{itemize}[itemsep=2pt]
  \item \texttt{--num-mcmc 0} disables the per-tempering-level filter
        HMC entirely. The HMC loop in \texttt{run\_outer\_smc}
        iterates 0 times.
  \item \texttt{--liu-west-a 0.97} applies per-dimension Liu--West
        shrinkage to the resampled $\theta$-cloud at every tempering
        level. (Skipped here in favour of \texttt{smooth-resample}.)
  \item \texttt{--smooth-resample-bw 1.0} switches the
        per-tempering-level rejuvenation to the multivariate
        Silverman-bandwidth KDE resample with coupled Liu--West
        correction $a = \sqrt{1 - h_{\mathrm{norm}}^2}$, ported from
        \texttt{SMC2FC\_functional/src/Filtering/Kernels.jl::}
        \texttt{smooth\_resample}. When set $>0$, this branch wins
        over the per-dim Liu--West.
  \item \texttt{--gaussian-bridge true} replaces the identity-copy of
        the previous-window posterior with a fresh sample from
        $\mathcal{N}(\hat{\mu}, \hat{\Sigma}_{\mathrm{reg}})$ via
        \texttt{Bridge.fit\_gaussian} +
        \texttt{Bridge.sample\_from\_gaussian}
        (regularised sample covariance, $10^{-6}$ Tikhonov).
  \item \texttt{--ctrl-n-smc 1024 --ctrl-num-mcmc 8} \textbf{REDUCES
        the controller particle / MCMC budget from the saturated
        default}: \texttt{--ctrl-n-smc} from $2048 \to 1024$ and
        \texttt{--ctrl-num-mcmc} from $16 \to 8$. This is the
        technical guide's ``shorter-wall-time''
        recipe (\texttt{compare\_v15\_julia\_vs\_python/docs/}
        \texttt{technical\_guide\_to\_current\_best\_Julia\_SMC2FC\_config.tex},
        ``When to deviate from these defaults''), traded against a
        weaker controller schedule. The technical guide quotes this
        change as approximately a $4\times$ wall-time reduction.
        \textbf{Controller HMC stays ON} at this reduced config —
        that code path is \texttt{run\_tempered\_smc\_gpu} in
        \texttt{Control/GPUControlSMC.jl}, separate from the filter
        HMC; only the filter HMC is being disabled.
  \item \texttt{--collect-ctrl-diagnostics true} writes one row of
        controller-HMC telemetry per tempering level to
        \texttt{controller\_diagnostics.csv} (new instrumentation,
        \S\ref{sec:ctrl_hmc}).
\end{itemize}

\subsection{Hardware}

NVIDIA GeForce RTX 5090; per-horizon \texttt{nvidia-smi} sampled at
1\,Hz to \texttt{T<N>d\_seed42/nvidia\_smi.csv}.


% =====================================================================
\section{Wall-time and closed-loop metrics}
\label{sec:wall}

Source: \texttt{horizons\_results.csv} (one row per horizon, written
by \texttt{compare\_v15\_julia\_vs\_python/src/extract\_horizon\_results.py}).

\begin{table}[ht]
\centering
\small
\begin{tabular}{rrrrrrr}
\toprule
$T$ (d) & wall (s) & $\overline{A}_{\mathrm{MPC}}$ & $A^{\mathrm{end}}_{\mathrm{MPC}}$ & $\overline{A}_{\mathrm{base}}$ & $A^{\mathrm{end}}_{\mathrm{base}}$ & $F$ violations \\
\midrule
14 & 133.2  & 0.077 & 0.085 & 0.114 & 0.182 & 0 \\
28 & 342.0  & 0.095 & 0.161 & 0.181 & 0.342 & 0 \\
42 & 715.5  & 0.167 & 0.521 & 0.285 & 0.646 & 0 \\
56 & 1234.6 & 0.332 & 1.066 & 0.401 & 0.824 & 0 \\
70 & 1790.8 & 0.491 & 1.197 & 0.492 & 0.894 & 0 \\
84 & 2456.6 & 0.612 & 1.238 & 0.565 & 0.907 & 0 \\
\bottomrule
\end{tabular}
\caption{Per-horizon wall time, mean and final autonomic amplitude
$A$ under MPC and under the constant-$\Phi=1$ baseline rollout, and
fatigue-barrier ($F_{\max}=0.40$) violation count. Zero $F$
violations across all six horizons.}
\label{tab:wall}
\end{table}

\subsection{GPU utilisation and memory}

\begin{table}[ht]
\centering
\small
\begin{tabular}{rrrr}
\toprule
$T$ (d) & mean GPU util (\%) & peak VRAM (MiB) & mean power (W) \\
\midrule
14 & 29.6 & 5399 & 181.0 \\
28 & 47.5 & 5390 & 280.3 \\
42 & 58.3 & 6091 & 335.7 \\
56 & 65.2 & 5273 & 369.7 \\
70 & 71.0 & 5272 & 398.5 \\
84 & 74.5 & 5776 & 415.4 \\
\bottomrule
\end{tabular}
\caption{Per-horizon GPU telemetry. Peak VRAM stays $\lesssim 6.1$\,GiB
(versus the 32\,GiB on the RTX 5090).}
\label{tab:gpu}
\end{table}

\subsection{Wall comparison vs the saturated-controller baseline}

The original 6-horizon baseline at
\texttt{julia\_horizon\_sweep\_2026-05-09/} ran the same six
horizons at the \emph{saturated} controller config
(\texttt{--ctrl-n-smc 2048 --ctrl-num-mcmc 16}) with the filter HMC
nominally on (but accepting 0\,\% of proposals). \textbf{This run
uses a reduced controller config} (\texttt{--ctrl-n-smc 1024
--ctrl-num-mcmc 8}, half the particles, half the MCMC moves). The
wall ratio below conflates two changes: removing the filter HMC and
halving the controller's particle / MCMC budget. The technical guide
attributes the bulk of the $\sim 4\times$ wall reduction to the
controller-side change alone:

\begin{table}[ht]
\centering
\small
\begin{tabular}{rrrr}
\toprule
$T$ (d) & this run (s) & baseline (s) & ratio \\
\midrule
14 & 133.2  & 422.7  & 3.17$\times$ \\
28 & 342.0  & 1259.3 & 3.68$\times$ \\
42 & 715.5  & 2286.6 & 3.20$\times$ \\
56 & 1234.6 & 3979.4 & 3.22$\times$ \\
70 & 1790.8 & 5586.4 & 3.12$\times$ \\
84 & 2456.6 & 7236.4 & 2.95$\times$ \\
\bottomrule
\end{tabular}
\caption{Wall comparison. \textbf{Both} the controller-config
reduction (\texttt{--ctrl-n-smc} $2048 \to 1024$,
\texttt{--ctrl-num-mcmc} $16 \to 8$) \textbf{and} the filter-HMC
removal contribute; the technical guide attributes the bulk of the
ratio to the controller-config change alone.}
\label{tab:wall_ratio}
\end{table}


% =====================================================================
\section{Controller-HMC measurements (new instrumentation)}
\label{sec:ctrl_hmc}

Each call to the controller's tempered-SMC$^2$ inner loop now records
per-tempering-level diagnostics (Listing~\ref{lst:ctrl_diag}). One
\texttt{controller\_diagnostics.csv} per horizon; aggregated by
\texttt{compare\_v15\_julia\_vs\_python/src/diag\_controller\_hmc.jl}
into per-horizon plots and a cross-horizon summary CSV.

\begin{lstlisting}[caption={Diagnostic columns recorded per tempering level.},
                    label={lst:ctrl_diag}, language=]
stride_idx, level, beta_pre, beta_post, delta_beta, beta_max,
ess_at_dbeta_chosen, n_smc, theta_dim,
chees_L_chosen, chees_score_best,
chees_L_candidates, chees_scores,
eps_step_size, n_mcmc_moves, total_accepts, accept_frac,
log_density_pre_mean, log_density_post_mean, delta_log_density,
esjd_total, esjd_per_eps_L, wall_seconds_level
\end{lstlisting}

\begin{table}[ht]
\centering
\footnotesize
\begin{tabular}{rrrrrrrrrr}
\toprule
$T$ & replans & total levels & mean accept (\%) & median accept (\%) & mean $L$ & mean $\beta_{\max}$ & mean ESJD/$\varepsilon L$ & mean $\Delta\log D$ & mean wall/level (s) \\
\midrule
14 & 13 & 69  & 97.93 & 97.91 & 16 & 158.20 & 14.08 & 0.012 & 1.69 \\
28 & 27 & 136 & 97.43 & 97.74 & 16 &  19.96 & 14.03 & 0.079 & 2.37 \\
42 & 41 & 230 & 96.17 & 96.48 & 16 &   7.16 & 13.41 & 0.344 & 3.00 \\
56 & 55 & 328 & 94.97 & 95.56 & 16 &   3.15 & 13.30 & 0.781 & 3.67 \\
70 & 69 & 412 & 93.49 & 94.78 & 16 &   1.82 & 13.20 & 1.194 & 4.27 \\
84 & 83 & 492 & 92.68 & 94.35 & 16 &   1.27 & 13.15 & 1.609 & 4.92 \\
\bottomrule
\end{tabular}
\caption{Cross-horizon controller-HMC summary (from
\texttt{docs/controller\_hmc\_summary.csv}). $L$ is the
ChEES-picked leapfrog count from the candidate list
$[16, 32, 64, 128, 256, 512]$; in this run the picker selects $L=16$
at every tempering level of every replan of every horizon.}
\label{tab:ctrl_hmc}
\end{table}

\paragraph{Filter-HMC accept-rate reference.} With \texttt{--num-mcmc 0}
the filter HMC loop iterates 0 times per tempering level; the
\texttt{accept=0\%} lines printed in the bench log are arithmetic
artefacts ($0 / \max(1, 0\cdot N_{\mathrm{smc}}) = 0$), \emph{not}
measurements of HMC quality. By contrast: in the original baseline
run (\texttt{julia\_horizon\_sweep\_2026-05-09/}) with the filter HMC
nominally on, 2880 of 2880 logged tempering levels reported 0\,\%
acceptance — that \emph{was} a measurement of attempted HMC moves.


% =====================================================================
\section{Posterior parameter traces (per horizon)}
\label{sec:param_traces}

Figures~\ref{fig:param_T14}--\ref{fig:param_T84}: per-horizon
posterior parameter trace plots, all written by the bench's internal
\texttt{\_plot\_param\_traces\_v15} function. Each figure is the same
10-panel layout: posterior median (blue), 5--95\,\% credible band
(light blue fill), truth value (red dashed), $x$-axis in days.

\begin{figure}[ht]
\centering
\includegraphics[width=\linewidth]{../T14d_seed42/v15_T14d_param_traces.png}
\caption{$T = 14$ days. Posterior parameter traces.}
\label{fig:param_T14}
\end{figure}

\begin{figure}[ht]
\centering
\includegraphics[width=\linewidth]{../T28d_seed42/v15_T28d_param_traces.png}
\caption{$T = 28$ days. Posterior parameter traces.}
\label{fig:param_T28}
\end{figure}

\begin{figure}[ht]
\centering
\includegraphics[width=\linewidth]{../T42d_seed42/v15_T42d_param_traces.png}
\caption{$T = 42$ days. Posterior parameter traces.}
\label{fig:param_T42}
\end{figure}

\begin{figure}[ht]
\centering
\includegraphics[width=\linewidth]{../T56d_seed42/v15_T56d_param_traces.png}
\caption{$T = 56$ days. Posterior parameter traces.}
\label{fig:param_T56}
\end{figure}

\begin{figure}[ht]
\centering
\includegraphics[width=\linewidth]{../T70d_seed42/v15_T70d_param_traces.png}
\caption{$T = 70$ days. Posterior parameter traces.}
\label{fig:param_T70}
\end{figure}

\begin{figure}[ht]
\centering
\includegraphics[width=\linewidth]{../T84d_seed42/v15_T84d_param_traces.png}
\caption{$T = 84$ days. Posterior parameter traces.}
\label{fig:param_T84}
\end{figure}

\clearpage


% =====================================================================
\section{Other plots on disk (not embedded — for separate inspection)}
\label{sec:other_plots}

\begin{itemize}[itemsep=2pt]
  \item \textbf{Cloud std (log-space) per stride, all 6 horizons
        overlaid:}
        \texttt{docs/cloud\_std\_per\_stride\_all\_horizons.png}.
  \item \textbf{Per-horizon controller-HMC 6-panel diagnostics:}
        \texttt{docs/controller\_hmc\_diagnostics\_T<N>d.png} for
        $N \in \{14, 28, 42, 56, 70, 84\}$.
  \item \textbf{Per-horizon state traces} ($B$, $F$, $A$, $\Phi$):
        \texttt{T<N>d\_seed42/v15\_T<N>d\_traces.png}.
\end{itemize}


% =====================================================================
\section{Source edits applied}
\label{sec:edits}

Two source files modified plus one new launcher and one new
diagnostic plotter:

\begin{itemize}[itemsep=2pt]
  \item \texttt{julia/SMC2FC\_functional/src/Control/GPUControlSMC.jl}
        — added \texttt{collect\_diagnostics::Bool=false} keyword on
        \texttt{run\_tempered\_smc\_gpu} (4-tuple return when true);
        added \texttt{return\_scores::Bool=false} keyword on
        \texttt{chees\_pick\_L\_generic}. Backward-compatible: existing
        callers with the default keyword get the original 3-tuple /
        2-tuple return.
  \item \texttt{version\_1\_5\_Julia/tools/bench\_smc\_full\_mpc\_fsa\_gpu.jl}
        — added \texttt{--collect-ctrl-diagnostics} CLI flag (default
        \texttt{true}), plumbed through \texttt{controller\_plan} and
        \texttt{stride\_step}; added \texttt{controller\_diagnostics.csv}
        writer. The manifest now records \texttt{liu\_west\_a},
        \texttt{smooth\_resample\_bw}, \texttt{gaussian\_bridge}.
  \item (NEW) \texttt{version\_1\_5\_Julia/tools/launchers/}\\
        \texttt{run\_julia\_horizon\_sweep\_no\_filter\_hmc.sh}
        — sibling sweep launcher with the new flag set inlined.
  \item (NEW) \texttt{compare\_v15\_julia\_vs\_python/src/}\\
        \texttt{diag\_controller\_hmc.jl} —
        aggregator + plotter for \texttt{controller\_diagnostics.csv}.
\end{itemize}


% =====================================================================
\section{Conclusion}
\label{sec:conclusion}

\textbf{Even for this very simple linear (additive Gaussian)
observation model, the filter \emph{without} the HMC kernel cannot
exactly recover the true model parameters.} Inspection of
Figures~\ref{fig:param_T14}--\ref{fig:param_T84} shows the posterior
median (blue) sitting at a consistent offset from truth (red dashed)
on several parameters at every horizon — the offset does not shrink
as $T$ grows.

This is a baseline. It says: rejuvenation tricks alone (Liu--West,
Silverman-KDE, Gaussian-bridge) are not sufficient to drive the
posterior to truth. Whatever is driving the residual bias — the
gradient information that an MCMC kernel would inject, the
between-window propagation choice, or something else — is the
question follow-up studies need to answer.


% =====================================================================
\section{What this report does NOT claim}
\label{sec:not_claimed}

\begin{itemize}[itemsep=2pt]
  \item Numerical per-parameter posterior medians extracted from
        \texttt{data.jld2} for each horizon. Not done — those would
        require loading
        \texttt{posterior\_particles[end\_stride, :, :]}, taking
        \texttt{log}, and computing the median, separately for each
        of the 6 horizons (and for the baseline run, if a
        side-by-side table is wanted). The per-parameter \emph{visual}
        offset is what \S\ref{sec:conclusion} relies on.
  \item Numerical cloud-std-vs-time curves. Visible in
        \texttt{docs/cloud\_std\_per\_stride\_all\_horizons.png}; not
        extracted as numbers in this report.
  \item Any causal attribution between which knob (\texttt{liu\_west\_a},
        \texttt{smooth\_resample\_bw}, \texttt{gaussian\_bridge}) is
        responsible for which observed effect. No ablations were run.
  \item Any verdict on the cause of the residual bias in
        \S\ref{sec:conclusion} — only the existence of the bias.
\end{itemize}


\end{document}
